% !TeX root = ../../Book3.tex
% 纲要标题：### 5.1 整元素
% 纲要位置：工作记录/计划纲要.md，第 1935 行。
% 纲要细目（写作范围）：
% * 首一方程及有限模判据
% * 整性的传递性
% * 整闭包及正规整环

\section{整元素}\label{b3:sec:0501}

有理数 \(1/2\) 满足整数系数方程 \(2T-1=0\)，却不满足整数系数的首一方程。首项系数能否取为一，决定了高次幂能否用低次幂表示，也决定了加入一个元素以后是否仍能用有限个模生成元控制运算。本节始终约定环交换且有单位，环包含保持单位。

\subsection{首一方程与有限模判据}\label{b3:subsec:0501:01}

\begin{definition}\label{b3:def:integral-element}
设 \(A\subseteq B\) 为环包含。若 \(b\in B\) 满足
\[
b^n+a_1b^{n-1}+\cdots+a_n=0,\qquad a_i\in A,
\]
则称 \(b\) 在 \(A\) 上\term[zheng-yuansu]{整}[integral]。若 \(B\) 的每个元素都在 \(A\) 上整，称 \(B/A\) 为\term[zheng-kuozhang]{整扩张}[integral extension]。若 \(B\) 作为 \(A\) 模有限生成，称 \(B/A\) 为\term[youxian-huan-kuozhang]{有限扩张}[finite extension]。
\end{definition}

这里的“有限”是模的有限生成性，不是底集合有限，也不是仅指代数生成元有限。例如 \(A[T]\) 是有限型 \(A\) 代数，但当 \(A\ne0\) 时不可能是有限 \(A\) 模。

\begin{proposition}[整性的有限模判据]\label{b3:prop:integral-finite-criterion}
对 \(b\in B\)，下列条件等价：
\begin{enumerate}[label=\Roman*.]
\item \(b\) 在 \(A\) 上整；
\item \(A[b]\) 是有限 \(A\) 模；
\item 存在有限 \(A\) 模 \(M\)，它同时是忠实 \(A[b]\) 模，且两种 \(A\) 作用一致。
\end{enumerate}
因此，有限环扩张必为整扩张；若 \(B=A[b_1,\ldots,b_r]\)，则 \(B/A\) 有限当且仅当每个 \(b_i\) 在 \(A\) 上整。
\end{proposition}
\begin{proof}
首一方程可把 \(b^n\) 及更高幂递次化为 \(1,b,\ldots,b^{n-1}\) 的 \(A\) 线性组合，故第一项推出第二项。取 \(M=A[b]\)，其中一使乘法作用忠实，得到第三项。

设第三项成立，选 \(A\) 模生成元 \(m_1,\ldots,m_s\)，写 \(bm_j=\sum_i c_{ij}m_i\)，其中 \(c_{ij}\in A\)。\cref{b3:thm:determinant-trick}说明首一多项式 \(\det(TI_s-C)\) 在乘法算子 \(b\) 上取值为零。忠实性再说明这个取值作为 \(A[b]\) 元素也为零，所以 \(b\) 整。

对有限扩张取 \(M=B\)，乘法作用同样忠实，便得到每个元素整。反过来，若 \(b_i\) 满足次数 \(n_i\) 的首一方程，则所有单项式
\[
b_1^{e_1}\cdots b_r^{e_r},\qquad 0\le e_i<n_i,
\]
生成整个 \(B\)，因为可依次降低每个变量的指数。
\end{proof}

例如 \(\mathbb Z[\sqrt2,\sqrt3]\) 由 \(1,\sqrt2,\sqrt3,\sqrt6\) 作为整数模生成。每个有限代数扩域都是整扩张；一般代数扩域也整，却未必有限。例如 \(\overline{\mathbb F}_p/\mathbb F_p\) 包含任意大的有限子域，不是有限维扩张。

\subsection{整性的传递性}\label{b3:subsec:0501:02}

\begin{proposition}[整性的传递性]\label{b3:prop:integral-transitivity}
若 \(A\subseteq B\subseteq C\)，\(B/A\) 和 \(C/B\) 都整，则 \(C/A\) 整。
\end{proposition}
\begin{proof}
对 \(c\in C\)，取其在 \(B\) 上的首一方程，设其中的有限个系数为 \(b_1,\ldots,b_r\)。由\cref{b3:prop:integral-finite-criterion}，\(B_0=A[b_1,\ldots,b_r]\) 是有限 \(A\) 模，\(B_0[c]\) 是有限 \(B_0\) 模。若两层分别由 \(u_i\) 与 \(v_j\) 生成，则乘积 \(u_iv_j\) 作为 \(A\) 模生成 \(B_0[c]\)。再次应用有限模判据，得到 \(c\) 在 \(A\) 上整。
\end{proof}

\begin{proposition}\label{b3:prop:integral-quotient-localization}
设 \(B/A\) 整。对理想 \(J\subseteq B\)，包含 \(A/(J\cap A)\subseteq B/J\) 整；对 \(A\) 中的乘法集 \(S\)，包含 \(S^{-1}A\subseteq S^{-1}B\) 整。
\end{proposition}
\begin{proof}
商环中的首一方程由原方程逐项取像得到。局部化中，对 \(b/s\) 把 \(b\) 的次数为 \(n\) 的首一方程除以 \(s^n\)，得到系数 \(a_i/s^i\in S^{-1}A\) 的首一方程。局部化的包含仍是单射：若 \(a/s\) 在 \(S^{-1}B\) 中为零，则某个 \(t\in S\) 满足 \(ta=0\)，该等式也在 \(A\) 中成立。
\end{proof}

这一命题允许把整性问题移到某个素理想附近，或移到某个闭集上。它不声称任意 \(B\) 中分母的局部化仍在 \(A\) 上整；例如把整数中的所有非零元素倒置便得到非整扩张 \(\mathbb Z\subseteq\mathbb Q\)。

\subsection{整闭包与正规整环}\label{b3:subsec:0501:03}

\begin{proposition}\label{b3:prop:integral-closure-ring}
\(B\) 中所有在 \(A\) 上整的元素组成一个包含 \(A\) 的子环 \(\overline A^{\,B}\)。若 \(c\in B\) 在此子环上整，则 \(c\in\overline A^{\,B}\)。
\end{proposition}
\begin{proof}
若 \(x,y\) 整，则 \(A[x,y]\) 是有限 \(A\) 模，其中的 \(x+y\)、\(xy\)、\(-x\) 均由\cref{b3:prop:integral-finite-criterion}在 \(A\) 上整。这给出子环性。最后一项由\cref{b3:prop:integral-transitivity}得到。
\end{proof}

\begin{definition}\label{b3:def:integral-closure}
子环 \(\overline A^{\,B}\) 称为 \(A\) 在 \(B\) 中的\term[zhengbi-bao]{整闭包}[integral closure]。若它等于 \(A\)，称 \(A\) 在 \(B\) 中整闭。整环 \(A\) 若在其分式域 \(\operatorname{Frac}(A)\) 中整闭，称为\term[zhengbi-zhenghuan]{整闭整环}[integrally closed domain]，也称\term[zhenggui-zhenghuan]{正规整环}[normal domain]。
\end{definition}

本章的正规整环不把 Noether 性包含在定义中。整闭性必须带有环境：\(\mathbb Z\) 在 \(\mathbb Q\) 中整闭，却不在 \(\mathbb Q(\sqrt2)\) 中整闭。

\begin{proposition}\label{b3:prop:ufd-integrally-closed}
唯一分解整环是整闭整环。整闭整环关于不含零的乘法集的局部化仍是整闭整环。
\end{proposition}
\begin{proof}
设 \(A\) 是唯一分解整环，\(a/b\) 为互素分子分母表示。若它满足次数 \(n\) 的首一方程，乘以 \(b^n\) 得 \(b\mid a^n\)。唯一分解及互素性迫使 \(b\) 为单位，故 \(a/b\in A\)。

设 \(A\) 整闭、\(0\notin S\)，\(z\in\operatorname{Frac}(A)\) 在 \(S^{-1}A\) 上整。把其方程的有限个系数通分为 \(a_i/s\)，其中 \(s\in S\)。令 \(w=sz\)，乘原方程以 \(s^n\)，得
\[
w^n+a_1w^{n-1}+a_2s w^{n-2}+\cdots+a_ns^{n-1}=0.
\]
于是 \(w\in A\)，从而 \(z=w/s\in S^{-1}A\)。
\end{proof}

特别地，域上的多项式环是整闭整环。相反，\(A=k[t^2,t^3]\subseteq k[t]\) 的分式域含有 \(t=t^3/t^2\)，而 \(t\) 满足 \(T^2-t^2=0\) 却不在 \(A\) 中，因为 \(A\) 的每个多项式的一次项系数为零。这是后面正规化计算的基本例子。
